% preamble.tex — shared LaTeX setup for Elementary Category Theory

\usepackage[T1]{fontenc}
\usepackage[utf8]{inputenc}
\usepackage{lmodern}
\usepackage{microtype}

% Page geometry
\usepackage[
  paperwidth=6in,
  paperheight=9in,
  top=0.85in,
  bottom=0.85in,
  inner=0.75in,
  outer=0.65in,
  marginparwidth=0in,
]{geometry}

% Typography
\usepackage{setspace}
\setstretch{1.15}
\usepackage{parskip}
\setlength{\parskip}{0.6em}
\setlength{\parindent}{0pt}

% Colors
\usepackage[dvipsnames,svgnames]{xcolor}
\definecolor{CoreColor}{HTML}{1A3A5C}
\definecolor{Exp1Color}{HTML}{1D6B3A}
\definecolor{Exp2Color}{HTML}{7B3F00}
\definecolor{Exp3Color}{HTML}{4A0072}
\definecolor{DialogBg}{HTML}{F5F0E8}
\definecolor{DialogBorder}{HTML}{C8A86B}
\definecolor{NoteBg}{HTML}{EEF4FF}
\definecolor{NoteBorder}{HTML}{6B8EC8}
\definecolor{TermColor}{HTML}{1A3A5C}

% Math
\usepackage{amsmath,amssymb,amsthm}
\usepackage{mathtools}

% Diagrams
\usepackage{tikz}
\usepackage{tikz-cd}
\usetikzlibrary{arrows,arrows.meta,positioning,calc,decorations.pathmorphing}
\tikzcdset{
  arrow style=math font,
  diagrams={>=stealth}
}

% Boxes and environments
\usepackage{tcolorbox}
\tcbuselibrary{skins,breakable,theorems}

% Section heading colors
\usepackage{titlesec}
\titleformat{\chapter}[display]
  {\normalfont\Large\bfseries\color{CoreColor}}
  {\chaptertitlename\ \thechapter}{12pt}{\Huge}
\titleformat{\section}
  {\normalfont\large\bfseries\color{CoreColor}}
  {\thesection}{1em}{}
\titleformat{\subsection}
  {\normalfont\normalsize\bfseries\color{Exp1Color}}
  {\thesubsection}{1em}{}
\titleformat{\subsubsection}
  {\normalfont\normalsize\bfseries\color{Exp2Color}}
  {\thesubsubsection}{1em}{}

% A fourth level for Expansion 3
\newcounter{deepexpand}[subsubsection]
\newcommand{\deepexpandsection}[1]{%
  \refstepcounter{deepexpand}%
  \vspace{0.8em}%
  {\noindent\normalfont\normalsize\bfseries\color{Exp3Color}%
    Deep Expansion \thesubsubsection.\thedeepexpand\quad #1}%
  \vspace{0.3em}\par%
}

% ── Expansion-level labels ──────────────────────────────────────────────────
% These appear as colored tags before expansion sections to make the
% staggered structure immediately visible to the reader.

\newcommand{\CoreLabel}{%
  \noindent\colorbox{CoreColor}{\color{white}\footnotesize\bfseries\strut\ CORE\ }%
  \vspace{0.3em}\par}

\newcommand{\ExpOneLabel}{%
  \noindent\colorbox{Exp1Color}{\color{white}\footnotesize\bfseries\strut\ EXPANSION 1\ }%
  \vspace{0.3em}\par}

\newcommand{\ExpTwoLabel}{%
  \noindent\colorbox{Exp2Color}{\color{white}\footnotesize\bfseries\strut\ EXPANSION 2\ }%
  \vspace{0.3em}\par}

\newcommand{\ExpThreeLabel}{%
  \noindent\colorbox{Exp3Color}{\color{white}\footnotesize\bfseries\strut\ EXPANSION 3\ }%
  \vspace{0.3em}\par}

% ── Internal-dialog box ─────────────────────────────────────────────────────
% Represents the annotated internal monologue a reader should develop
% while working through diagrams and definitions.

\newtcolorbox{dialog}[1][]{
  enhanced,
  breakable,
  colback=DialogBg,
  colframe=DialogBorder,
  fonttitle=\bfseries\small,
  title={Internal Dialog},
  left=6pt, right=6pt, top=4pt, bottom=4pt,
  attach boxed title to top left={yshift=-2mm, xshift=4mm},
  boxed title style={colback=DialogBorder, colframe=DialogBorder,
    fontupper=\color{white}\bfseries\small},
  #1
}

% ── Margin note (reader annotation reminder) ────────────────────────────────
\newtcolorbox{readernote}[1][]{
  enhanced,
  breakable,
  colback=NoteBg,
  colframe=NoteBorder,
  fonttitle=\bfseries\small,
  title={Reader's Annotation},
  left=6pt, right=6pt, top=4pt, bottom=4pt,
  attach boxed title to top left={yshift=-2mm, xshift=4mm},
  boxed title style={colback=NoteBorder, colframe=NoteBorder,
    fontupper=\color{white}\bfseries\small},
  #1
}

% ── Key-term command ────────────────────────────────────────────────────────
\newcommand{\term}[1]{\textbf{\color{TermColor}#1}}

% ── Example environment ─────────────────────────────────────────────────────
\newtcolorbox{example}[2][]{
  enhanced,
  breakable,
  colback=white,
  colframe=CoreColor!40,
  fonttitle=\bfseries\small,
  title={Example: #2},
  left=6pt, right=6pt, top=4pt, bottom=4pt,
  #1
}

% ── Summary box ─────────────────────────────────────────────────────────────
\newtcolorbox{summarybox}[1][]{
  enhanced,
  breakable,
  colback=CoreColor!5,
  colframe=CoreColor,
  fonttitle=\bfseries,
  title={Chapter Summary},
  left=8pt, right=8pt, top=6pt, bottom=6pt,
  #1
}

% Hyperlinks (load last)
\usepackage[hidelinks,
  pdfauthor={},
  pdftitle={Elementary Category Theory}]{hyperref}
